% --- Archivo: 06_sqp.tex ---

% =========================================================
\section{Programación Cuadrática Secuencial (SQP)}
\label{sec:v2_metodos_sqp}
% =========================================================

Como su nombre indica, los métodos de \deftech{Programación Cuadrática Secuencial} (SQP, por sus siglas en inglés: \textit{Sequential Quadratic Programming}) consisten en la resolución de una sucesión de problemas de programación cuadrática, que aproximan en cada iteración la topología local del problema no lineal subyacente. 

Un problema de programación cuadrática (minimizar una función cuadrática bajo restricciones afines) puede ser resuelto eficientemente de manera analítica o con algoritmos de complejidad polinómica finita. Esta estrategia es el generalizador natural de la idea del método de Newton clásico para problemas con restricciones.

\subsection{El subproblema SQP general}

Dada una aproximación del iterado primal $x^k$, y sus multiplicadores duales asociados $\lambda^k$ y $\mu^k$, la función objetivo se modela como una aproximación de Taylor de segundo orden y las restricciones (igualdad $h$ y desigualdad $g$) se relajan mediante expansiones afines de primer orden.

El subproblema de programación cuadrática estándar para determinar la dirección de búsqueda $d \in \Rn$ se formula explícitamente como:
\begin{equation}
    \min_{d \in \Rn} \interno{f'(x^k)}{d} + \frac{1}{2}\interno{H_k d}{d}
    \label{eq:v2_sqp_obj}
\end{equation}
sujeto a las restricciones linealizadas que definen el poliedro de avance $D_k$:
\begin{align}
    h(x^k) + J_h(x^k)d &= 0, \label{eq:v2_sqp_rest_h} \\
    g(x^k) + J_g(x^k)d &\le 0. \label{eq:v2_sqp_rest_g}
\end{align}

El bloque analítico crítico recae sobre la construcción y determinación de la matriz simétrica de curvatura $H_k \in \R^{n \times n}$. Un análisis puramente geométrico podría sugerir usar $H_k = \nabla^2 f(x^k)$. Sin embargo, dicha elección omite completamente la aceleración centrípeta y la curvatura no lineal que dictan las fronteras del conjunto factible $D$.

La justificación Newtoniana formal exige que la curvatura global considere simultáneamente la topología de la función y de las restricciones, siendo la matriz de diseño óptimo asintótica:
\begin{equation}
    H_k = \nabla_{xx}^2 L(x^k, \lambda^k, \mu^k),
    \label{eq:v2_sqp_hessiana_exacta}
\end{equation}
donde $L$ es la función Lagrangiana clásica.

\begin{algorithm}[El método SQP general]\label{alg:v2_sqp_general}
    Elegir un punto de partida $(x^0, \lambda^0, \mu^0) \in \Rn \times \R^l \times \R_+^m$ y una matriz de curvatura $H_0$. Tomar $k := 0$.
    \begin{enuthm}
        \item Calcular la dirección óptima $d^k \in \Rn$ resolviendo el programa cuadrático exacto dictado por \eqref{eq:v2_sqp_obj}-\eqref{eq:v2_sqp_rest_g}. Se obtienen de manera paralela y simultánea los nuevos multiplicadores de Lagrange óptimos $(y^k, z^k)$ asociados a las restricciones linealizadas del subproblema.
        \item Tomar el paso completo de Newton sobre el espacio primal y reemplazar los multiplicadores duales:
        \[
            x^{k+1} = x^k + d^k, \quad \lambda^{k+1} = y^k, \quad \mu^{k+1} = z^k.
        \]
        \item Actualizar la matriz $H_{k+1}$ evaluando la nueva Hessiana de la Lagrangiana (o aplicando una actualización Cuasi-Newton tipo BFGS amortiguada). Incrementar $k \leftarrow k + 1$ y regresar al Paso 1.
    \end{enuthm}
\end{algorithm}

\begin{theorem}[Convergencia local superlineal del método SQP]\label{thm:v2_sqp_local}
    Supongamos que $\bar{x}$ es un punto estacionario del problema original que cumple estrictamente la Independencia Lineal de Gradientes Activos (LICQ), la condición Suficiente de Segundo Orden, y exhibe Complementariedad Estricta ($\bar{\mu}_i > 0 \ \forall i \in I(\bar{x})$).
    
    Si se utiliza la Hessiana de la Lagrangiana exacta \eqref{eq:v2_sqp_hessiana_exacta}, entonces para cualquier punto de partida $(x^0, \lambda^0, \mu^0)$ suficientemente cercano a la solución analítica, el \cref{alg:v2_sqp_general} está bien definido (los subproblemas no colapsan y admiten solución) y converge a la solución óptima $(\bar{x}, \bar{\lambda}, \bar{\mu})$ con \textbf{tasa de convergencia cuadrática}.
\end{theorem}

\subsection{Estrategias de Globalización del Método SQP}

A grandes rasgos, el esquema SQP base del \cref{alg:v2_sqp_general} adolece de las mismas fallas analíticas que el método de Newton puro sin regularizar: no tiene garantía de convergencia si se inicializa fuera de la vecindad de atracción del óptimo, y el subproblema \eqref{eq:v2_sqp_obj} puede divergir si la matriz $H_k$ asume autovalores negativos (o si las restricciones linealizadas resultan ser mutuamente incompatibles, haciendo que $D_k = \emptyset$).

Para dotar al algoritmo SQP de garantías de convergencia global, se implementa obligatoriamente un proceso de búsqueda lineal (Line Search) o de regiones de confianza. Esto se logra acoplando y guiando el descenso mediante una \textbf{Función de Mérito}.

\subsubsection{La Función de Mérito Penalizada Exacta}

Dado que en optimización restringida debemos equilibrar la reducción de $f(x)$ con la eliminación progresiva de la violación topológica de la factibilidad, introducimos una función de mérito escalar exacta $L_1$:
\begin{equation}
    \varphi(x; c) = f(x) + c \left( \norm{h(x)}_1 + \sum_{i=1}^m \max\{0, g_i(x)\} \right),
    \label{eq:v2_merito_L1_sqp}
\end{equation}
donde el peso penalizador $c > 0$ es un parámetro dinámico que debe exceder el límite superior de los multiplicadores (es decir, $c > \norm{(\lambda, \mu)}_\infty$). 

Se puede demostrar rigurosamente que si la matriz de curvatura $H_k$ es simétrica definida positiva, la dirección $d^k$ producida por el subproblema SQP es consistentemente una \textbf{dirección de descenso garantizada} para la función de mérito $\varphi$. Esto autoriza la implementación segura de la regla de Armijo y evita la in-factibilidad general.

\subsubsection{El Efecto Maratos y las fallas de la tasa superlineal}

La convergencia superlineal rápida del método SQP requiere asintóticamente que la búsqueda lineal permita tomar pasos puros completos (es decir, $\alpha_k = 1$). Sin embargo, existe una patología topológica conocida como el \deftech{Efecto Maratos}.

En problemas reales, si el iterado $x^k$ se ubica peligrosamente cerca de la frontera factible curvada, dar un paso completo tangente $d^k$ puede resultar topológicamente en un leve alejamiento (incremento residual) de la frontera factible curva. Este minúsculo incremento residual penalizado impacta desproporcionadamente en la función de mérito exacta $L_1$, causando que $\varphi(x^k + d^k; c) > \varphi(x^k; c)$.

En respuesta, el algoritmo rechaza iterativamente el paso completo (exigiendo $\alpha_k < 1$), destruyendo estrepitosamente la tasa cuadrática de convergencia y obligando al método a avanzar con extrema lentitud. Este efecto perturbador y su mitigación (utilizando pasos de corrección de segundo orden transversales o relajando condicionalmente las funciones de mérito) es un tema crucial y altamente activo en el diseño práctico de rutinas comerciales tipo \textit{fmincon} en el estado de arte del análisis matemático no lineal.